\section{结论：稀疏容量的收益来自选择，代价也来自选择}

主讲结论可以压缩为三个判断。第一，SMoE 用稀疏 activation 换取更大知识容量，并能处在优秀的 active-compute frontier；第二，训练和 serving 能做到高效，但绝不只由 active parameter count 决定；第三，expert specialization 有结构，却不像“代码专家/数学专家”那样直接。最后一张 slide 把 architecture 与 interpretability research 同时保留为开放方向。

\lecturefigure{slide-26-conclusion.jpg}{Albert Jiang 对 Sparse MoE、效率与 expert specialization 的总结}{Stanford Online 官方录像 00:29:38}

\begin{importantbox}{本讲最终公式：三本账必须一起报告}
\[
\boxed{\text{Model decision}=f(P_{\text{total}},P_{\text{active}},M_{\text{resident}},C_{\text{comm}},I_{\text{load}},Q_{\text{quality}})}
\]
其中 total/active 描述结构，resident memory/communication/load 描述系统，quality 描述任务证据。任何只保留一个变量的比较都可能误导。
\end{importantbox}

\teachervoice{00:29:18--00:29:47，讲者的原始总结是：SMoE 借 sparsity 获得更多 knowledge，可以训练成高效 inference model，但 expert specialization 没有直觉上那么简单，architecture 与 interpretability 仍有大量工作。}

\subsection{本章小结}

Mixtral 的价值不在“八个专家”这个标签，而在 token-level conditional computation：总容量增长、活跃路径受控、router 动态选择。它同时创造了新的系统瓶颈和新的可解释观测量。只有把结构、硬件和证据放在一起，才能判断某个 MoE 是否真的优于 dense baseline。

